\chapter{Über FreeQDA}
FreeQDA ist ein freies open-source Softwareprojekt zur Analyse qualitativer Forschungsdaten.


\begin{table}[!h]
\centering
\begin{tabular}{p{30mm}p{40mm}|p{60mm}}
%\textbf{trallala} & \textbf{juppida} \\
\multirow{5}{*}{\includegraphics[width=25mm]{img/freeQDA_FINAL_1024}} & \textbf{Programmierung} & Dirk Kitscha	\\
&\textbf{Konzept} & Jörg große Schlarmann	\\
&\textbf{Beratung} & André Fringer		\\
&\textbf{Homepage} & \url{http://freeqda.sf.net} \\
&\textbf{Lizenz} & \href{http://www.gnu.org/licenses/gpl-2.0.html}{GPL 2.0}		\\
\end{tabular}
%\caption{Über FreeQDA}
\label{tab:about}
\end{table}
\vfill
FreeQDA ist freie Software und darf unter den Bedingungen der \href{http://www.gnu.org/licenses/gpl-2.0.html}{GPL 2.0} %
frei verwendet, kopiert und weiterentwickelt werden.\\
Die Entwicklung von FreeQDA wurde von der \textit{Deutschen Gesellschaft für Pflegewissenschaft e.V.}\footnote{\url{http://www.dg-pflegewissenschaft.de}} %
(DGP) gefördert.
\vfill